\subsection{注记与进一步结果}
\label{sec:ch8-notes-further-results}

\subsubsection{参考文献}
\label{subsec:ch8-references}

J. E. Littlewood 和 R. A. E. C. Paley 的原始结果见 \textit{Theorems on Fourier series and power series, I} (J. London Math. Soc. 6 (1931), 230--233) 和 \textit{Theorems on Fourier series and power series, II} (Proc. London Math. Soc. 42 (1936), 52--89). 他们使用了复变量方法, 另见 Zygmund [21]. 本章采用的方法源自 A. Benedek, A. P. Calderon 和 R. Panzone, \textit{Convolution operators with Banach-space valued functions}, Proc. Nat. Acad. Sci. U.S.A. 48 (1962), 356--365. 另见 J. L. Rubio de Francia, F. J. Ruiz 和 J. L. Torrea, \textit{Calderón--Zygmund theory for operator-valued kernels}, Adv. in Math. 62 (1986), 7--48.

Hörmander 乘子定理最早见 \textit{Estimates for translation invariant operators in $L^p$-spaces}, Acta Math. 104 (1960), 93--139. Marcinkiewicz 乘子定理见 \textit{Sur les multiplicateurs des séries de Fourier}, Studia Math. 8 (1939), 78--91. 正文所给证明不同于原始证明. 在 Hörmander 定理中还可证明乘子为弱 $(1,1)$ 型. Bochner--Riesz 乘子由 S. Bochner 在 \textit{Summation of multiple Fourier series by spherical means}, Trans. Amer. Math. Soc. 40 (1936), 175--207 中引入. 第~\ref{sec:ch8-return-singular-integrals}~节和第~\ref{sec:ch8-parabola-maximal-hilbert}~节的方法改编自 J. Duoandikoetxea 和 J. L. Rubio de Francia, \textit{Maximal and singular integral operators via Fourier transform estimates}, Invent. Math. 84 (1986), 541--561. 关于曲线上的奇异积分与极大函数, 见下文第~\ref{subsec:ch8-operators-curves}~节. Van der Corput 引理最早见 \textit{Zahlentheoretische Abschätzungen}, Math. Ann. 84 (1921), 53--79.

Littlewood--Paley 理论还有其他版本, 包括使用不同平方函数或连续参数的版本, 例如见 Stein [15, Chapter 4]. 定理~\ref{thm:ch8-smooth-square-rn} 可看作借助 Littlewood--Paley 分解 $\{S_jf\}$ 对 $L^p$ 空间的刻画. 类似地, 许多其他空间, 包括 Besov 空间和 Triebel--Lizorkin 空间, 都可用 Littlewood--Paley 分解的不同混合范数不等式来刻画. 进一步内容见 M. Frazier, B. Jawerth 和 G. Weiss, \textit{Littlewood--Paley theory and the study of function spaces}, CBMS Regional Conference Series in Mathematics 79, Amer. Math. Soc., Providence, 1991, 以及 H. Triebel, \textit{Theory of function spaces II}, Birkhäuser, Basel, 1992.

\hypertarget{chapter-08-page-197}{}
\subsubsection{其他 Littlewood--Paley 不等式}
\label{subsec:ch8-other-littlewood-paley}

\begin{Theorem}
\label{thm:ch8-arbitrary-interval-square-function}
设 $\{I_j\}$ 是 $\mathbb R$ 中任意一列互不相交的区间, 并由
\[
 \widehat{S_jf}(\xi)=\chi_{I_j}(\xi)\widehat f(\xi)
\]
定义 $S_j$. 则当 $2\le p<\infty$ 时,
\[
 \left\|\left(\sum_j|S_jf|^2\right)^{1/2}\right\|_p
 \le C_p\|f\|_p.
\]
\end{Theorem}

当 $I_j=(j,j+1)$ 时, 该定理最早由 L. Carleson 在 \textit{On the Littlewood--Paley theorem}, Report, Mittag-Leffler Inst., Djursholm, 1967 中证明. 后来的新证明见 A. Córdoba, \textit{Some remarks on the Littlewood--Paley theory}, Rend. Circ. Mat. Palermo 1 (1981), suppl., 75--80, 以及 J. L. Rubio de Francia, \textit{Estimates for some square functions of Littlewood--Paley type}, Publ. Mat. 27 (1983), 81--108. 一般结果由 J. L. Rubio de Francia 在 \textit{A Littlewood--Paley inequality for arbitrary intervals}, Rev. Mat. Iberoamericana 1 (1985), 1--14 中证明, 另一证明由 J. Bourgain 在 \textit{On square functions on the trigonometric system}, Bull. Soc. Math. Belg. Sér. B 37 (1985), 20--26 中给出. 取 $\widehat f_N=\chi_{[0,N]}$ 即可看出, 当 $1<p<2$ 时该结论不成立.

定理~\ref{thm:ch8-arbitrary-interval-square-function} 在 $\mathbb R^n$ 中的推广由 J. L. Journé 在 \textit{Calderón--Zygmund operators on product spaces}, Rev. Mat. Iberoamericana 1 (1985), 55--91 中证明. 与 $n=1$ 的证明相比, 其证明相当困难. 后来较简单的证明由 P. Sjölin, \textit{A note on Littlewood--Paley decompositions with arbitrary intervals}, J. Approx. Theory 48 (1986), 328--334, F. Soria, \textit{A note on a Littlewood--Paley inequality for arbitrary intervals in $\mathbb R^2$}, J. London Math. Soc. 36 (1987), 137--142, 以及 S. Sato, \textit{Note on a Littlewood--Paley operator in higher dimensions}, J. London Math. Soc. 42 (1990), 527--534 给出.

定理~\ref{thm:ch8-arbitrary-interval-square-function} 可用来改进 Marcinkiewicz 乘子定理~\ref{thm:ch8-marcinkiewicz-multiplier}.

\hypertarget{chapter-08-page-198}{}
\begin{Theorem}
\label{thm:ch8-square-bounded-variation-multiplier}
若 $m$ 有界, 且 $m^2$ 在 $\mathbb R$ 的每个二进区间上有界变差, 则 $m$ 是 $L^p$, $1<p<\infty$, 上的乘子.
\end{Theorem}

该结果由 R. Coifman, J. L. Rubio de Francia 和 S. Semmes 在 \textit{Multiplicateurs de Fourier de $L^p(\mathbb R)$ et estimations quadratiques}, C. R. Acad. Sci. Paris 306 (1988), 351--354 中证明.

在 $\mathbb R^2$ 中还可证明一个角向 Littlewood--Paley 定理. 令 $\Delta_j$ 为斜率分别为 $2^j$ 和 $2^{j+1}$ 的两条直线所围成的区域, 并由 $\widehat{S_jf}(\xi)=\chi_{\Delta_j}(\xi)\widehat f(\xi)$ 定义 $S_j$. 则对 $1<p<\infty$,
\[
 \left\|\left(\sum_j|S_jf|^2\right)^{1/2}\right\|_p
 \le C_p\|f\|_p.
\]
该结果由 A. Nagel, E. M. Stein 和 S. Wainger 在 \textit{Differentiation in lacunary directions}, Proc. Nat. Acad. Sci. U.S.A. 75 (1978), 1060--1062 中证明.

Littlewood--Paley 不等式也有加权版本. 当 $w\in A_p$ 时, 定理~\ref{thm:ch8-smooth-square-rn} 在 $L^p(w)$ 中成立. D. Kurtz 在 \textit{Littlewood--Paley and multiplier theorems on weighted $L^p$ spaces}, Trans. Amer. Math. Soc. 259 (1980), 235--254 中对支于环带的 $\psi$ 证明了这一点, 其结论可推广至定理~\ref{thm:ch8-smooth-square-rn} 中的 $\psi$. Rubio de Francia 对定理~\ref{thm:ch8-arbitrary-interval-square-function} 的证明实际上在 $p>2$ 时给出带 $A_{p/2}$ 权的加权版本, 其证明基于他与 Ruiz 和 Torrea 的上述论文所引入的方法. 高维中的同一结果由 H. Lin 在 \textit{Some weighted inequalities on product domains}, Trans. Amer. Math. Soc. 318 (1990), 69--85 中证明. 尚不知道定理~\ref{thm:ch8-arbitrary-interval-square-function} 的加权版本对 $L^2(w)$, $w\in A_1$, 是否成立. 但对等距区间这一特殊情形, 该结论成立, 见上引 Rubio de Francia 在 Publ. Mat. 发表的论文.

\hypertarget{chapter-08-page-199}{}
\subsubsection{球的乘子与 Bochner--Riesz 乘子}
\label{subsec:ch8-ball-bochner-riesz}

在维数 $n\ge2$ 时, 球的乘子表现出最坏的可能性质.

\begin{Theorem}
\label{thm:ch8-ball-multiplier-negative}
当 $n\ge2$ 且 $p\ne2$ 时, $\mathbb R^n$ 中 Euclidean 球的示性函数不是 $L^p$ 乘子.
\end{Theorem}

这一否定结果由 C. Fefferman 在 \textit{The multiplier problem for the ball}, Ann. of Math. 94 (1972), 330--336 中证明. 它远强于由引理~\ref{lem:ch8-compact-multiplier-kernel-lp} 和引理~\ref{lem:ch8-bochner-riesz-kernel} 所能得到的结论, 因为后一方法只能在 $|1/p-1/2|\ge1/(2n)$ 时证明不是乘子. 不过, 限制到 $L^p$ 中的径向函数时, 若 $|1/p-1/2|<1/(2n)$, 它是乘子. 这一结果由 C. S. Herz 在 \textit{On the mean inversion of Fourier and Hankel transforms}, Proc. Nat. Acad. Sci. U.S.A. 40 (1954), 996--999 中证明.

对 $a>0$, 定理~\ref{thm:ch8-bochner-riesz} 中的充分条件最早由 E. M. Stein 借助定理~\ref{thm:ch1-riesz-thorin} 在 \textit{Interpolation of linear operators}, Trans. Amer. Math. Soc. 83 (1956), 482--492 中证明. 当 $0<a<(n-1)/2$ 时, 该结果遗漏了一段 $p$ 值范围. 尽管 $a=0$, 即球的乘子, 有上述否定结果, 仍猜想当
\begin{equation}
 \left|\frac1p-\frac12\right|<\frac{2a+1}{2n}
 \label{eq:ch8-bochner-riesz-conjectured-range}
\end{equation}
时, $T^a$ 在 $L^p$ 上有界.

该猜想只在 $\mathbb R^2$ 中得到完全证明. 第一个证明由 L. Carleson 和 P. Sjölin 在 \textit{Oscillatory integrals and a multiplier problem for the disc}, Studia Math. 44 (1972), 287--299 中给出. 其他证明见 L. Hörmander, \textit{Oscillatory integrals and multipliers on $FL^p$}, Ark. Mat. 11 (1973), 1--11, C. Fefferman, \textit{A note on spherical summation multipliers}, Israel J. Math. 15 (1973), 44--52, 以及 A. Córdoba, \textit{A note on Bochner--Riesz operators}, Duke Math. J. 46 (1979), 505--511.

在高维中, 首个对某些 $a<(n-1)/2$ 证明~\eqref{eq:ch8-bochner-riesz-conjectured-range} 所给全部 $p$ 值范围的结果由 C. Fefferman 在 \textit{Inequalities for strongly singular convolution operators}, Acta Math. 124 (1970), 9--36 中得到. 他证明了 $a>(n-1)/4$ 的情形. 其证明的一个重要方面是引入了 Bochner--Riesz 乘子与 Fourier 变换到球面的限制之间的联系. 下一节给出由此联系导出的更强结果. 这些结果均可见 K. M. Davis 和 Y. C. Chang [3].

\hypertarget{chapter-08-page-200}{}
Bochner--Riesz 乘子给出 Fourier 变换的一族求和方法. 除范数收敛外, 还可考虑逐点收敛, 即是否
\begin{equation}
 \lim_{R\to\infty}T_R^af(x)=f(x)\quad\text{a.e.},
 \label{eq:ch8-bochner-riesz-pointwise-convergence}
\end{equation}
其中
\[
 \widehat{T_R^af}(\xi)=\left(1-\frac{|\xi|^2}{R^2}\right)_+^a\widehat f(\xi).
\]
研究逐点收敛的通常方法是考察相应的极大算子 $(T^a)^*f(x)=\sup_{R>0}|T_R^af(x)|$, 见第~\ref{sec:ch2-hl-maximal-function}~节. 当 $a$ 大于临界指标 $(n-1)/2$ 时, $(T^a)^*$ 由 Hardy--Littlewood 极大算子控制, 因而逐点收敛立即成立. 当 $a<(n-1)/2$ 时已有部分结果. 若 $n=2$ 且 $p>2$, 则在~\eqref{eq:ch8-bochner-riesz-conjectured-range} 所给范围内~\eqref{eq:ch8-bochner-riesz-pointwise-convergence} 成立, 见 A. Carbery, \textit{The boundedness of the maximal Bochner--Riesz operator on $L^4(\mathbb R^2)$}, Duke Math. J. 50 (1983), 409--416. 对 $n\ge3$, $p>2$ 且 $a>(n-1)/(2(n+1))$ 的同一结果由 M. Christ 在 \textit{On almost everywhere convergence of Bochner--Riesz means in higher dimensions}, Proc. Amer. Math. Soc. 95 (1985), 16--20 中证明. 两个结果都通过证明对应 $p$ 值下 $(T^a)^*$ 在 $L^p$ 上有界得到.

A. Carbery, J. L. Rubio de Francia 和 L. Vega 在 \textit{Almost everywhere summability of Fourier integrals}, J. London Math. Soc. 38 (1988), 513--524 中证明, 若 $0<\beta<1+2a<n$, 则 $(T^a)^*$ 在加权空间 $L^2(|x|^{-\beta})$ 上有界, 并推出这些空间中的几乎处处收敛. 对满足~\eqref{eq:ch8-bochner-riesz-conjectured-range} 的 $p>2$, 存在 $\beta<1+2a$ 使 $L^p\subset L^2+L^2(|x|^{-\beta})$, 因此~\eqref{eq:ch8-bochner-riesz-pointwise-convergence} 对所有这样的 $p$ 成立. 这给出的 Bochner--Riesz 平均几乎处处收敛的 $p$ 值范围, 大于目前已知 $T^a$, 更不用说 $(T^a)^*$, 有界的范围.

\subsubsection{限制定理}
\label{subsec:ch8-restriction-theorems}

若 $f\in L^1$, 则 $\widehat f$ 连续且有界, 因而对任意 $S\subset\mathbb R^n$,
\[
 \|\widehat f\|_{L^\infty(S)}\le\|f\|_{L^1(\mathbb R^n)}.
\]
但一般而言, 若 $f\in L^p(\mathbb R^n)$, $p>1$, 且 $S$ 的测度为零, 则 $\widehat f$ 仅几乎处处有定义, 因而不能先验地定义其在 $S$ 上的值. 如果对 $f\in C_c^\infty$ 有估计
\[
 \|\widehat f\|_{L^q(S)}\le C_{p,q}\|f\|_{L^p(\mathbb R^n)},
\]
则由连续性, 可把 $f\in L^p$ 的 Fourier 变换在 $S$ 上的限制定义为 $L^q$ 函数. 若 $S$ 包含于某个超平面, 该估计仅在 $p=1$ 时成立. 令人惊讶的是, 对某些非零曲率的 $S$ 可以得到非平凡限制估计.

\hypertarget{chapter-08-page-201}{}
单位球面 $S=S^{n-1}$ 的情形尤其重要, 因为它与 Bochner--Riesz 乘子有关. C. Fefferman 的下述定理说明了这一点, 见前引 Acta 论文.

\begin{Theorem}
\label{thm:ch8-restriction-implies-bochner-riesz}
若对某个 $p>1$ 有
\begin{equation}
 \|\widehat f\|_{L^2(S^{n-1})}
 \le C_p\|f\|_{L^p(\mathbb R^n)},
 \label{eq:ch8-sphere-l2-restriction}
\end{equation}
则对满足~\eqref{eq:ch8-bochner-riesz-conjectured-range} 的所有 $a$, $T^a$ 在 $L^p$ 上有界.
\end{Theorem}

C. Fefferman 对 $1<p<4n/(3n+1)$ 证明了限制不等式, 这给出上文提到的条件 $a>(n-1)/4$. P. Tomas 在发展 E. M. Stein 早期工作的基础上改进了该结果, 见 \textit{A restriction theorem for the Fourier transform}, Bull. Amer. Math. Soc. 81 (1975), 477--478.

\begin{Theorem}[Tomas--Stein]
\label{thm:ch8-tomas-stein}
对所有满足
\[
 1\le p\le\frac{2n+2}{n+3}
\]
的 $p$, 不等式~\eqref{eq:ch8-sphere-l2-restriction} 成立.
\end{Theorem}

直接推论是, 若 $a>(n-1)/(2n+2)$, 则 Bochner--Riesz 乘子在~\eqref{eq:ch8-bochner-riesz-conjectured-range} 的全部范围内有界.

一般猜想
\begin{equation}
 \|\widehat f\|_{L^q(S^{n-1})}
 \le C_{p,q}\|f\|_{L^p(\mathbb R^n)}
 \label{eq:ch8-general-restriction}
\end{equation}
成立当且仅当
\[
 1\le p<\frac{2n}{n+1},
 \qquad
 \frac1p>\frac{n+1}{n-1}\frac1q.
\]
$p,q$ 必须限制在此范围内, 这可通过取 $f$ 为一个适配球面的长方体的示性函数看出. 这里的长方体以 $(1,0,\ldots,0)$ 为中心, 边长为 $\delta\times\delta^{1/2}\times\cdots\times\delta^{1/2}$. 当 $n=2$ 时, C. Fefferman 在前引 Acta 论文中证明了该猜想. 另见 A. Zygmund, \textit{On Fourier coefficients and transforms of functions of two variables}, Studia Math. 50 (1974), 189--201. 当 $q=2$ 时猜想也成立, 这由定理~\ref{thm:ch8-tomas-stein} 给出.

在 $n\ge3$ 时, 定理~\ref{thm:ch8-tomas-stein} 长期是最佳已知结果, 直到 J. Bourgain 在 \textit{Besicovitch type maximal operators and applications to Fourier analysis}, Geom. Funct. Anal. 1 (1991), 147--187 中加以改进. 他证明了限制猜想在 $1<p<p(n)$ 时成立, 其中 $p(n)$ 递归定义并满足
\[
 \frac{2n+2}{n+3}<p(n)<\frac{2n}{n+1}.
\]
例如 $p(3)=31/23$.

\hypertarget{chapter-08-page-202}{}
限制问题、Bochner--Riesz 乘子与 Kakeya 极大算子之间联系密切, 见第~\ref{subsec:ch2-kakeya-maximal}~节. Bourgain 的论文改进了这些问题此前的全部已知结果.

Bochner--Riesz 猜想和限制问题的进一步改进仍由 J. Bourgain 给出, 见 \textit{Some new estimates on oscillatory integrals}, Essays on Fourier Analysis in Honor of Elias M. Stein, C. Fefferman, R. Fefferman and S. Wainger, eds., pp. 83--112, Princeton Univ. Press, Princeton, 1995. 当 $n=3$ 时, 这些结果又由 T. Tao, A. Vargas 和 L. Vega 在 \textit{A bilinear approach to the restriction and Kakeya conjectures}, J. Amer. Math. Soc. 11 (1998), 967--1000 中改进. T. Tao 还证明了定理~\ref{thm:ch8-ball-multiplier-negative} 的一种逆命题, 见 \textit{The Bochner--Riesz conjecture implies the restriction conjecture}, Duke Math. J. 96 (1999), 363--376.

\subsubsection{乘子的加权范数不等式}
\label{subsec:ch8-weighted-multipliers}

给定乘子 $m$ 和相应算子 $T_m$, 可以像第~\ref{chap:weighted-inequalities}~章处理奇异积分那样, 尝试刻画使 $T_m$ 在 $L^p(w)$ 上有界的权 $w$. Marcinkiewicz 乘子定理可直接推广到加权情形.

\begin{Theorem}
\label{thm:ch8-weighted-marcinkiewicz}
设 $m$ 是有界函数, 并且在 $\mathbb R$ 的每个二进区间上一致有界变差. 若 $w\in A_p$, 则 $T_m$ 在 $L^p(\mathbb R,w)$, $1<p<\infty$, 上有界.
\end{Theorem}

该结果由 D. Kurtz 证明, 见第~\ref{subsec:ch8-other-littlewood-paley}~节所引论文. 他还在该文中证明了定理~\ref{thm:ch8-marcinkiewicz-multiplier-r2} 的加权对应版本.

对 Hörmander 定理~\ref{cor:ch8-hormander-classical} 中的乘子, $A_p$ 条件不再充分, 但有以下结论.

\begin{Theorem}
\label{thm:ch8-weighted-hormander-multiplier}
设 $n\ge2$, $k=[n/2]+1$. 若 $n/k<p<\infty$ 且 $w\in A_{pk/n}$, 或 $1<p<(n/k)'$ 且 $w^{1-p'}\in A_{p'k/n}$, 则对满足~\eqref{eq:ch8-hormander-annular-condition} 的乘子 $m$, 算子 $T_m$ 在 $L^p(w)$ 上有界. 若 $n>2$, 则 $p$ 的范围还可包含端点 $p=n/k$ 和 $p=(n/k)'$.
\end{Theorem}

定理~\ref{thm:ch8-weighted-hormander-multiplier} 由 D. Kurtz 和 R. Wheeden 在 \textit{Results on weighted norm inequalities for multipliers}, Trans. Amer. Math. Soc. 255 (1979), 343--362 中证明.

Bochner--Riesz 乘子的加权范数不等式所知少得多. 当 $a\ge(n-1)/2$, $w\in A_p$ 时, 已知 $T^a$ 在 $L^p(w)$, $1<p<\infty$, 上有界. 当 $a>(n-1)/2$ 时可证明
\[
 |T^af(x)|\le CMf(x),
\]
其中 $M$ 是 Hardy--Littlewood 极大算子, 因而结论立即由定理~\ref{thm:ch7-maximal-weighted-strong} 得到. 对临界指标 $a=(n-1)/2$, 该结果由 X. Shi 和 Q. Sun 在 \textit{Weighted norm inequalities for Bochner--Riesz operators and singular integral operators}, Proc. Amer. Math. Soc. 116 (1992), 665--673 中证明. 另一证明可由上引 Duoandikoetxea 和 Rubio de Francia 的论文得到. 在此情形, 算子对 $A_1$ 权还是弱 $(1,1)$ 型, 见 A. Vargas, \textit{Weighted weak type $(1,1)$ bounds for rough operators}, J. London Math. Soc. 54 (1996), 297--310. 一般情形的部分结果见 K. Andersen, \textit{Weighted norm inequalities for Bochner--Riesz spherical summation multipliers}, Proc. Amer. Math. Soc. 103 (1988), 165--171.

\hypertarget{chapter-08-page-203}{}
\subsubsection{球面极大函数}
\label{subsec:ch8-spherical-maximal}

推论~\ref{cor:ch8-dyadic-spherical-maximal} 证明了当 $n\ge2$ 时, 二进球面极大函数在 $L^p$, $1<p<\infty$, 上有界. 这一结论最早由 C. Calderon, \textit{Lacunary spherical means}, Illinois J. Math. 23 (1979), 476--484, 以及 R. Coifman 和 G. Weiss, \textit{Review of the book Littlewood--Paley and multiplier theory by R. E. Edwards and G. I. Gaudry}, Bull. Amer. Math. Soc. 84 (1978), 242--250 独立证明. 对连续参数极大函数
\[
 \mathcal Mf(x)=\sup_{t>0}\int_{S^{n-1}}f(x-ty)\,d\sigma(y),
\]
有以下结论.

\begin{Theorem}
\label{thm:ch8-spherical-maximal}
$\mathcal M$ 在 $L^p(\mathbb R^n)$ 上有界当且仅当 $p>n/(n-1)$.
\end{Theorem}

对 $n\ge3$, 定理~\ref{thm:ch8-spherical-maximal} 由 E. M. Stein 在 \textit{Maximal functions: Spherical means}, Proc. Nat. Acad. Sci. U.S.A. 73 (1976), 2174--2175 中证明. 对 $n=2$, 由 J. Bourgain 在 \textit{Averages in the plane over convex curves and maximal operators}, J. Analyse Math. 47 (1986), 69--85 中证明. $n\ge3$ 时的其他证明由 M. Cowling 和 G. Mauceri, \textit{On maximal functions}, Rend. Sem. Mat. Fis. Milano 49 (1979), 79--87, 以及 J. L. Rubio de Francia, \textit{Maximal functions and Fourier transforms}, Duke Math. J. 53 (1986), 395--404 给出. 后一个证明在思想上接近第~\ref{sec:ch8-return-singular-integrals}~节的方法, 即把乘子分解为二进部分, 对每部分作带指数衰减的 $L^2$ 估计, 再用向量值方法导出 $L^1$ 估计.

在 $n=2$ 时, 新证明见 G. Mockenhaupt, A. Seeger 和 C. Sogge, \textit{Wave front sets, local smoothing and Bourgain's circular maximal theorem}, Ann. of Math. 136 (1992), 207--218, 以及 W. Schlag, \textit{A geometric proof of the circular maximal theorem}, Duke Math. J. 93 (1998), 505--533. 球面极大算子在离散和连续参数下的加权范数不等式见 J. Duoandikoetxea 和 L. Vega, \textit{Spherical means and weighted inequalities}, J. London Math. Soc. 53 (1996), 343--353.

\hypertarget{chapter-08-page-204}{}
\subsubsection{奇异积分的进一步结果}
\label{subsec:ch8-further-singular-integrals}

第~\ref{sec:ch8-parabola-maximal-hilbert}~节的论证可直接推广, 以证明无法用旋转法得到的其他奇异积分结果. 例如, 定义 $T$ 为与核
\[
 K(x)=h(|x|)\frac{\Omega(x')}{|x|^n}
\]
的卷积, 其中对某个 $q>1$, $\Omega\in L^q(S^{n-1})$ 且积分为零, 而 $h$ 满足增长条件
\[
 \sup_{R>0}\frac1R\int_0^R|h(t)|^2\,dt<\infty.
\]
容易修改定理~\ref{thm:ch8-compact-kernel-dyadic-sum} 的证明, 得到 $T$ 在 $L^p$, $1<p<\infty$, 上有界. 此类奇异积分最早由 R. Fefferman, \textit{A note on singular integrals}, Proc. Amer. Math. Soc. 74 (1979), 266--270, 以及 J. Namazi, \textit{A singular integral}, Proc. Amer. Math. Soc. 96 (1986), 421--424 考虑.

类似方法也可用于证明加权范数不等式.

\begin{Theorem}
\label{thm:ch8-weighted-rough-singular-integral}
给定 $\Omega\in L^\infty(S^{n-1})$, 满足 $\int\Omega(u)\,d\sigma(u)=0$, 定义奇异积分
\[
 Tf(x)=\operatorname{p.v.}\int_{\mathbb R^n}
 \frac{\Omega(y')}{|y|^n}f(x-y)\,dy.
\]
若 $w\in A_p$, $1<p<\infty$, 则 $T$ 在 $L^p(w)$ 上有界.
\end{Theorem}

该结果的证明见上引 Duoandikoetxea 和 Rubio de Francia 的论文. 假设 $\Omega\in L^\infty$ 不能弱化为 $\Omega\in L^q$, $q<\infty$. 反例最早由 B. Muckenhoupt 和 R. Wheeden 在 \textit{Weighted norm inequalities for singular and fractional integrals}, Trans. Amer. Math. Soc. 161 (1971), 249--258 中给出. 当 $\Omega\in L^q$, $1<q<\infty$ 时, 若 $q'\le p<\infty$ 且 $w\in A_{p/q'}$, 或 $1<p\le q$ 且 $w^{1-p'}\in A_{p'/q'}$, 则 $T$ 在 $L^p(w)$ 上有界, 比较定理~\ref{thm:ch8-weighted-hormander-multiplier}. 这一条件由 J. Duoandikoetxea, \textit{Weighted norm inequalities for homogeneous singular integrals}, Trans. Amer. Math. Soc. 336 (1993), 869--880, 和 D. Watson, \textit{Weighted estimates for singular integrals via Fourier transform estimates}, Duke Math. J. 60 (1990), 389--399 独立发现.

\hypertarget{chapter-08-page-205}{}
这些算子的其他加权不等式以及权的结构性质, 见 D. Watson, \textit{Vector-valued inequalities, factorization, and extrapolation for a family of rough operators}, J. Funct. Anal. 121 (1994), 389--415, 以及 J. Duoandikoetxea, \textit{Almost-orthogonality and weighted inequalities}, Contemp. Math. 189 (1995), 213--226. 加权弱 $(1,1)$ 不等式由 Vargas 在第~\ref{subsec:ch8-weighted-multipliers}~节所引论文中证明.

\subsubsection{曲线上的算子}
\label{subsec:ch8-operators-curves}

除第~\ref{sec:ch8-parabola-maximal-hilbert}~节所述动机例子外, 曲线上的奇异积分和极大函数还出现在多种背景中, 并已得到广泛研究. 最早的结果针对齐次曲线, 即 $(t^{a_1},t^{a_2},\ldots,t^{a_n})$ 型曲线及其某些更一般版本. 后来这些结果推广到良好弯曲曲线, 即其导数 $\Gamma'(0),\Gamma''(0),\ldots$ 张成 $\mathbb R^n$ 的曲线 $\Gamma$. E. M. Stein 和 S. Wainger 的综述 \textit{Problems in harmonic analysis related to curvature}, Bull. Amer. Math. Soc. 84 (1978), 1239--1295 包含主要结果及其证明. Wainger 在 [16] 中的 \textit{Averages and singular integrals over lower dimensional sets} 也是很好的概述.

随后研究的是平坦曲线, 即在原点的所有导数都为零的曲线. 这类曲线首先在 $\mathbb R^2$ 中得到研究. 除上引 Duoandikoetxea 和 Rubio de Francia 的论文外, 还见 M. Christ, \textit{Hilbert transforms along curves, II: A flat case}, Duke Math. J. 52 (1985), 887--894; A. Nagel, J. Vance, S. Wainger 和 D. Weinberg, \textit{Hilbert transforms for convex curves}, Duke Math. J. 50 (1983), 735--744, 以及 \textit{Maximal function for convex curves}, Duke Math. J. 52 (1985), 715--722; H. Carlsson et al., \textit{$L^p$ estimates for maximal functions and Hilbert transforms along flat convex curves in $\mathbb R^2$}, Bull. Amer. Math. Soc. 14 (1986), 263--267; A. Córdoba 和 J. L. Rubio de Francia, \textit{Estimates for Wainger's singular integrals along curves}, Rev. Mat. Iberoamericana 2 (1986), 105--117; A. Carbery, M. Christ, J. Vance, S. Wainger 和 D. Watson, \textit{Operators associated to flat plane curves: $L^p$ estimates via dilation methods}, Duke Math. J. 59 (1989), 675--700.

平坦曲线结果向高维的推广更为困难. 参见 A. Nagel, J. Vance, S. Wainger 和 D. Weinberg, \textit{The Hilbert transform for convex curves in $\mathbb R^n$}, Amer. J. Math. 108 (1986), 485--504; A. Carbery, J. Vance, S. Wainger 和 D. Watson, \textit{The Hilbert transform and maximal function along flat curves, dilations, and differential equations}, Amer. J. Math. 116 (1994), 1203--1239; A. Carbery, J. Vance, S. Wainger, D. Watson 和 J. Wright, \textit{$L^p$ estimates for operators associated to flat curves without the Fourier transform}, Pacific J. Math. 167 (1995), 243--262; A. Carbery 和 S. Ziesler, \textit{Hilbert transforms and maximal functions along rough flat curves}, Rev. Mat. Iberoamericana 10 (1994), 379--393.

\hypertarget{chapter-08-page-206}{}
另见 A. Carbery, J. Vance, S. Wainger 和 D. Watson 的综述 \textit{Dilations associated to flat curves in $\mathbb R^n$}, Essays on Fourier Analysis in Honor of Elias M. Stein, C. Fefferman, R. Fefferman and S. Wainger, eds., pp. 113--126, Princeton Univ. Press, Princeton, 1995.

近年来还研究了非卷积算子. 此时算子定义在可变曲线上, 即曲线还依赖于点 $x$. 类似地, 在卷积情形和可变情形中都用曲面替代了曲线. 一些例子可见 Stein [17, Chapter 11].
